\section{\texorpdfstring{هفته‌ی 14: مدل‌سازی در علوم شناختی 1: دسته‌بندی
\lr{(GCM)} و تصمیم‌گیری
\lr{(Drift-Diffusion)}}{هفته‌ی 14: مدل‌سازی در علوم شناختی 1: دسته‌بندی (GCM) و تصمیم‌گیری (Drift-Diffusion)}}

\begin{quote}
\textbf{فصل 6: نمونه‌هایی از مدل‌سازی در حوزه‌های مختلف علوم شناختی} · درس
1 ساعت · \lr{lab} 2 ساعت \textbf{داده‌ی اصلی:} داده‌ی دسته‌بندی شبیه‌سازی‌شده
برای ساختار کلاسیک 5-4، و داده‌ی انتخاب\lr{--}\lr{RT} شبیه‌سازی‌شده از یک فرایند
\lr{drift-diffusion} در یک طرح سرعت/دقت × آسان/سخت.
\end{quote}

\begin{center}\rule{0.5\linewidth}{0.5pt}\end{center}

\subsection{واژه‌های کلیدی}

واژه‌های مهم این درس با معادل فارسی‌شان (حتی اگر در درس‌های قبل آمده
باشند).

\begin{longtable}[]{@{}
  >{\raggedright\arraybackslash}p{(\linewidth - 4\tabcolsep) * \real{0.3333}}
  >{\raggedright\arraybackslash}p{(\linewidth - 4\tabcolsep) * \real{0.3333}}
  >{\raggedright\arraybackslash}p{(\linewidth - 4\tabcolsep) * \real{0.3333}}@{}}
\caption{\centering واژه‌های کلیدی}\tabularnewline
\toprule\noalign{}
\begin{minipage}[b]{\linewidth}\raggedright
\lr{English}
\end{minipage} & \begin{minipage}[b]{\linewidth}\raggedright
فارسی
\end{minipage} & \begin{minipage}[b]{\linewidth}\raggedright
توضیح کوتاه
\end{minipage} \\
\midrule\noalign{}
\endfirsthead
\toprule\noalign{}
\begin{minipage}[b]{\linewidth}\raggedright
\lr{English}
\end{minipage} & \begin{minipage}[b]{\linewidth}\raggedright
فارسی
\end{minipage} & \begin{minipage}[b]{\linewidth}\raggedright
توضیح کوتاه
\end{minipage} \\
\midrule\noalign{}
\endhead
\bottomrule\noalign{}
\endlastfoot
\lr{categorization} & دسته‌بندی & قرار دادن محرک در یک دسته \\
\lr{Generalized Context Model (GCM)} & مدل بافت تعمیم‌یافته & مدل
\lr{exemplar} دسته‌بندی بر پایه‌ی شباهت \\
\lr{drift rate} & نرخ رانش & سرعت میانگین جمع شدن شواهد؛ نشانه‌ی کیفیت
شواهد \\
\lr{simulation} & شبیه‌سازی & تولید داده‌ی مصنوعی با اجرای مدل \\
\lr{response time (RT)} & زمان واکنش & فاصله‌ی ارائه‌ی محرک تا پاسخ \\
\lr{drift-diffusion model (DDM)} & مدل \lr{drift-diffusion}
(رانش\lr{--}انتشار) & انباشت پیوسته‌ی شواهد بین دو مرز \\
\lr{speed\lr{--}accuracy trade-off} & بده‌بستان سرعت و دقت & سریع‌تر پاسخ دادن
به قیمت خطای بیشتر \\
\lr{accuracy} & دقت & نسبت پاسخ‌های درست \\
\lr{exemplar} & نمونه \lr{(exemplar)} & یک مورد ذخیره‌شده در حافظه از یک
دسته \\
\lr{prototype} & پیش‌نمونه \lr{(prototype)} & خلاصه یا میانگین یک دسته \\
\lr{fit} / \lr{fitting} & برازش & تنظیم پارامترها تا مدل با داده جور
شود \\
\lr{attention weight} & وزن توجه & سهم توجه به هر بُعد محرک \\
\lr{sensitivity (d′)} & حساسیت & فاصله‌ی توزیع علامت و نویز \\
\lr{parameter} & پارامتر & کمیت قابل تنظیم مدل که بین افراد یا شرایط فرق
می‌کند \\
\lr{estimate} / \lr{estimator} & برآورد / برآوردگر & مقدار به‌دست‌آمده
برای پارامتر / روش به دست آوردن آن \\
\lr{EZ-diffusion} & \lr{EZ-diffusion} & برآورد ساده‌ی پارامترهای \lr{DDM}
از سه آماره \\
\lr{parameter recovery} & بازیابی پارامتر & آیا برازش، پارامترهای معلوم
شبیه‌سازی را برمی‌گرداند \\
\lr{evidence} & شواهد & اطلاعاتی که برای تصمیم جمع می‌شود \\
\lr{stimulus} & محرک & چیزی که به شرکت‌کننده نشان داده می‌شود \\
\lr{boundary separation} & فاصله‌ی مرزها & مقدار شواهد لازم برای پاسخ؛
نشانه‌ی احتیاط \\
\lr{non-decision time} & زمان غیرتصمیمی & زمان دیدن محرک و اجرای حرکت
پاسخ \\
\lr{similarity} & شباهت & نزدیکی دو محرک در فضای روان‌شناختی \\
\lr{trial} & کوشش \lr{(trial)} & یک بار انجام تکلیف توسط شرکت‌کننده \\
\lr{embedding} & \lr{embedding} (بردار معنایی) & برداری که مدل برای یک
کلمه یا متن می‌سازد \\
\end{longtable}

\begin{center}\rule{0.5\linewidth}{0.5pt}\end{center}

\subsection{اهداف یادگیری}

در پایان این هفته، دانشجویان می‌توانند:

\begin{enumerate}
\def\labelenumi{\arabic{enumi}.}
\tightlist
\item
  نظریه‌های نمونه‌محور \lr{(exemplar)} و پیش‌نمونه‌ای \lr{(prototype)}
  دسته‌بندی را \textbf{توصیف کنند} و آن‌ها را به‌صورت \lr{GCM} و یک مدل
  \lr{prototype} \textbf{صورت‌بندی کنند}.
\item
  \lr{GCM} و مدل \lr{prototype} را با \lr{MLE} روی داده‌ی دسته‌بندی
  \textbf{برازش کنند}، \textbf{مقایسه کنند}، و وزن‌های توجه و حساسیت را
  \textbf{تفسیر کنند}.
\item
  مدل \lr{drift-diffusion (DDM)} را به‌عنوان حد پیوسته‌ی \lr{random walk}
  هفته‌ی 1 \textbf{توصیف کنند} و پارامترهایش را به سازه‌های شناختی
  \textbf{نگاشت کنند}.
\item
  پارامترهای \lr{DDM} را با روش بسته‌ی \lr{(closed-form) EZ-diffusion}
  \textbf{برآورد کنند} و آن را با بازیابی پارامتر \textbf{اعتبارسنجی
  کنند}.
\item
  اثرهای آزمایشی را بر حسب پارامترهای مدل \textbf{تفسیر کنند} («این
  دست‌کاری احتیاط را تغییر داد، نه کیفیت شواهد را»).
\end{enumerate}

\begin{center}\rule{0.5\linewidth}{0.5pt}\end{center}

\subsection{مقدمه}

13 هفته‌ی اول یک جعبه‌ابزار ساختند، بیشتر روی دو مثال کوچک (فراموشی و
\lr{SDT).} فصل 6 جعبه‌ابزار را روی مدل‌هایی به کار می‌گیرد که کل حوزه‌های
پژوهشی را شکل داده‌اند. این هفته دو تا از تأثیرگذارترین‌ها را می‌بینیم.

\textbf{دسته‌بندی.} چطور تصمیم می‌گیریم که یک حیوان تازه سگ است؟
\emph{نظریه‌ی \lr{prototype}} می‌گوید آن را با یک خلاصه‌ی دسته، یعنی سگ
میانگین، مقایسه می‌کنیم. \emph{نظریه‌ی \lr{exemplar}} می‌گوید آن را با
خاطره‌های ذخیره‌شده‌ی همه‌ی سگ‌ها (و همه‌ی گربه‌ها)یی که دیده‌ایم مقایسه می‌کنیم.
هر دو معقول به نظر می‌رسند و بحث‌های کلامی سال‌ها رفت و برگشت. مدل‌های فرمال
بودند که فرق را ساختند: ساختار دسته‌ی 5-4 مدین و شیفر طوری طراحی شد که دو
نظریه برای محرک‌های مشخصی \textbf{پیش‌بینی‌های متضاد} داشته باشند.
\textbf{مدل بافت تعمیم‌یافته \lr{(GCM)}} نوسوفسکی نظریه‌ی \lr{exemplar} را
به یک مدل دقیق و قابل برازش با وزن‌های توجه تفسیرپذیر تبدیل کرد.

\textbf{تصمیم‌گیری.} در هفته‌ی 1 یک \lr{random walk} ساختیم که شواهد را تا
رسیدن به یک مرز جمع می‌کند. نسخه‌ی پیوسته‌اش، \textbf{مدل
\lr{drift-diffusion}}، احتمالاً موفق‌ترین مدل روان‌شناسی شناختی است: دقت
\emph{و} کل توزیع‌های \lr{RT} پاسخ‌های درست و خطا را توضیح می‌دهد، با
پارامترهایی برای کیفیت شواهد، احتیاط و \lr{non-decision time.} برای
مطالعه‌ی پیری، \lr{ADHD}، افسردگی، یادگیری و اثر داروها از آن استفاده
می‌شود. برازش \lr{DDM} کامل از نظر محاسباتی سنگین است، ولی روش
\textbf{\lr{EZ-diffusion}} با معادله‌های بسته از سه آماره‌ی خلاصه
برآوردهای خوبی می‌دهد.

\paragraph{این موضوع کجا به کار می‌رود؟}

\begin{itemize}
\tightlist
\item
  \textbf{ایده‌ی \lr{exemplar} در یادگیری ماشین:} طبقه‌بندهای نزدیک‌ترین
  همسایه \lr{(nearest neighbor)} و روش‌های \lr{kernel} یک مورد جدید را بر
  اساس شباهتش به مثال‌های ذخیره‌شده طبقه‌بندی می‌کنند؛ همان اصل \lr{GCM.}
\item
  \textbf{مدل \lr{drift-diffusion}} یکی از پرکاربردترین مدل‌ها در علوم
  اعصاب شناختی و پژوهش بالینی است: سالمندی، \lr{ADHD}، افسردگی، اضطراب،
  اثر داروها و بی‌خوابی. در پژوهش مصرف‌کننده (اینکه آدم‌ها با چه سرعتی بین
  محصولات انتخاب می‌کنند) و در عصب‌اقتصاد هم استفاده می‌شود.
\item
  \textbf{\lr{EZ-diffusion}} محبوب است چون برای هر نفر فقط سه آماره‌ی
  خلاصه لازم دارد و به بهینه‌سازی نیاز ندارد، پس می‌شود آن را سریع روی
  داده‌های موجود اجرا کرد.
\end{itemize}

\begin{center}\rule{0.5\linewidth}{0.5pt}\end{center}

\Needspace{11\baselineskip}

\subsection{جزوه‌ی درس}

\paragraph{به زبان ساده}

این هفته دو مدل معروف داریم. \textbf{دسته‌بندی:} چطور تصمیم می‌گیریم یک
حیوان تازه سگ است؟ یک نظریه می‌گوید آن را با یک «سگ ایده‌آل» میانگین
مقایسه می‌کنیم (\textbf{\lr{prototype}}). \textbf{\lr{GCM}} ایده‌ی
\lr{exemplar} را به ریاضی تبدیل می‌کند: فاصله‌ی آیتم جدید با هر مثال
ذخیره‌شده را بسنج (با \textbf{توجه} بیشتر به بعضی ویژگی‌ها)، فاصله را به
شباهت تبدیل کن (نزدیک‌تر = خیلی شبیه‌تر)، و دسته‌ی \lr{A} را با احتمال
«شباهت کل به \lr{A} / شباهت کل به همه» انتخاب کن. یک مجموعه محرک
هوشمندانه (ساختار 5-4) کاری می‌کند که دو نظریه برای بعضی آیتم‌ها
پیش‌بینی‌های \emph{متضاد} داشته باشند. \textbf{تصمیم‌گیری:} \textbf{مدل
\lr{drift-diffusion}} نسخه‌ی پیوسته‌ی \lr{random walk} هفته‌ی 1 است. شواهد
با سرعت میانگین \(v\) (کیفیت شواهد) به سمت مرز درست می‌رود و هم‌زمان نویز
آن را تکان می‌دهد؛ وقتی به مرز \(a\) (احتیاط) برسد پاسخ داده می‌شود؛ یک
زمان ثابت \(T_{er}\) هم برای دیدن و زدن کلید اضافه می‌شود.
\textbf{\lr{EZ-diffusion}} یک میان‌بر است که \(v\)، \(a\) و \(T_{er}\) را
فقط از سه عدد حساب می‌کند: دقت، و میانگین و واریانس زمان واکنش پاسخ‌های
درست.

\textbf{زمان‌بندی پیشنهادی (60 دقیقه):} نظریه‌های دسته‌بندی و ساختار 5-4 10
· \lr{GCM} 15 · مدل \lr{prototype} و پیش‌بینی‌های کلیدی 5 · مدل
\lr{drift-diffusion} 15 · \lr{EZ-diffusion} 10 · تفسیر 5.

\subsubsection{1. ساختار دسته‌ی 5-4}

محرک‌ها روی چهار بعد دودویی فرق دارند (مثلاً شکل، رنگ، اندازه و موقعیت یک
شکل هندسی در طرح مدین و شیفر؛ ویژگی‌های چهره در نسخه‌های بعدی).
شرکت‌کننده‌ها یاد می‌گیرند 9 محرک آموزشی را در دسته‌های \lr{A} و \lr{B}
بگذارند، و بعد 7 محرک \textbf{انتقال} \lr{(transfer)} تازه را بدون
بازخورد دسته‌بندی می‌کنند. مجموعه محرک رایج این است {[}با \lr{Medin} \lr{\&}
\lr{Schaffer,} 1978 یا \lr{Smith} \lr{\&} \lr{Minda,} 2000 بررسی شود{]}:

\begin{longtable}[]{@{}lll@{}}
\caption{\centering 1. ساختار دسته‌ی 5-4}\tabularnewline
\toprule\noalign{}
دسته‌ی \lr{A} & دسته‌ی \lr{B} & انتقال \\
\midrule\noalign{}
\endfirsthead
\toprule\noalign{}
دسته‌ی \lr{A} & دسته‌ی \lr{B} & انتقال \\
\midrule\noalign{}
\endhead
\bottomrule\noalign{}
\endlastfoot
\lr{A1} 1110 & \lr{B1} 1100 & \lr{T1} 1001 \\
\lr{A2} 1010 & \lr{B2} 0110 & \lr{T2} 1000 \\
\lr{A3} 1011 & \lr{B3} 0001 & \lr{T3} 1111 \\
\lr{A4} 1101 & \lr{B4} 0000 & \lr{T4} 0010 \\
\lr{A5} 0111 & & \lr{T5} 0101 \\
& & \lr{T6} 0011 \\
& & \lr{T7} 0100 \\
\end{longtable}

\lr{prototype} دسته‌ی \lr{A} برابر 1111 و \lr{prototype} دسته‌ی \lr{B}
برابر 0000 است. مقایسه‌ی کلیدی: \textbf{\lr{A1} (1110)} سه ویژگی با
\lr{prototype} دسته‌ی \lr{A} مشترک دارد، ولی \textbf{\lr{A2} (1010)} فقط
دو تا. پس مدل \lr{prototype} پیش‌بینی می‌کند \lr{A1} \emph{بهتر} از
\lr{A2} دسته‌بندی شود. اما \lr{A2} به دو \lr{exemplar} دیگر \lr{A} (یعنی
\lr{A1} و \lr{A3)} خیلی شبیه است و به هیچ \lr{exemplar} از \lr{B} شبیه
نیست، در حالی که \lr{A1} به دو \lr{exemplar} از \textbf{\lr{B}} (یعنی
\lr{B1} و \lr{B2)} خیلی شبیه است. پس مدل \lr{exemplar} می‌تواند پیش‌بینی
کند که \lr{A2} \emph{دست‌کم به خوبیِ} \lr{A1} دسته‌بندی شود، و در داده‌های
تجربی \lr{A2} معمولاً بهتر دسته‌بندی می‌شود.

\subsubsection{\texorpdfstring{2. مدل بافت تعمیم‌یافته
\lr{(GCM)}}{2. مدل بافت تعمیم‌یافته (GCM)}}

\textbf{فاصله} بین محرک \(i\) و \lr{exemplar} ذخیره‌شده‌ی \(e\)، با وزن‌های
توجه \(w_m \ge 0\) و \(\sum_m w_m = 1\) (متریک \lr{city-block} برای
ابعاد جدایی‌پذیر):

\[
\mathrm{dist}_{ie} = \sum_{m=1}^{4} w_m\, |x_{im} - x_{em}|.
\]

\textbf{شباهت} با فاصله به‌صورت نمایی کم می‌شود، با \textbf{حساسیت} \(c\):

\[
s_{ie} = \exp(-c \cdot \mathrm{dist}_{ie}).
\]

\textbf{قاعده‌ی انتخاب} (شباهت جمع‌شده، نسبت \lr{Luce):}

\[
P(A \mid i) = \frac{\sum_{e \in A} s_{ie}}{\sum_{e \in A} s_{ie} + \sum_{e \in B} s_{ie}}.
\]

\textbf{پارامترها:} \(c\) (شباهت با فاصله چقدر تند کم می‌شود؛ \(c\) بزرگ
یعنی فقط \lr{exemplar}های خیلی نزدیک مهم‌اند) و سه وزن توجه آزاد (چهارمی
\lr{\$}1 - \lr{\$} بقیه است)، پس \(k = 4\). (این همان \(c\) چندمعنایی جدول
نمادگذاری هفته‌ی 1 است؛ این هفته یعنی حساسیت، نه ملاک \lr{SDT.)}

\textbf{تفسیر وزن‌های توجه:} \(w_m\) سهم توجهی است که به بعد \(m\) داده
می‌شود. توجه بهینه روی ابعادی تأکید می‌کند که برای دسته تشخیصی
\lr{(diagnostic)} هستند. در ساختار 5-4، ابعاد 1 و 3 تشخیصی‌ترین‌اند.

\textbf{مدل \lr{prototype:}} همان فاصله، شباهت و قاعده‌ی انتخاب، ولی هر
دسته فقط با \lr{prototype} خودش نمایش داده می‌شود:
\(P(A\mid i) = s_{i,P_A}/(s_{i,P_A} + s_{i,P_B})\). این هم \(k = 4\)
دارد، پس شمارش پارامترها روی مقایسه اثری ندارد.

\textbf{درست‌نمایی:} اگر محرک \(i\) به تعداد \(n_i\) بار دسته‌بندی شود و
\(y_i\) بار \lr{«A»} نامیده شود،
\(y_i \sim \mathrm{Binomial}(n_i,\allowbreak P(A\mid i))\) است و با \lr{MLE} و
\lr{multi-start} برازش می‌شود (هفته‌های 3\lr{--}4). وزن‌ها بهتر است با یک
\textbf{\lr{reparameterization softmax}} (هفته‌ی 3) برآورد شوند،
\(w = \mathrm{softmax}(\eta_1,\allowbreak \eta_2,\allowbreak \eta_3,\allowbreak 0)\)، و
\(c = e^{\gamma}\).

\begin{quote}
\textbf{پیشرفته:} \lr{GCM} کامل یک پارامتر سوگیری پاسخ \(b_A\)، یک توان
شباهت (شباهت گاوسی در برابر نمایی)، یک پارامتر متریک \(r\)
\lr{(city-block} با \(r = 1\) در برابر اقلیدسی با \(r = 2\)) و گاهی یک
پارامتر قطعیت پاسخ \lr{(response determinism)} \(\gamma\) (به توان
رساندن شباهت‌های جمع‌شده) اضافه می‌کند. \lr{Lee} \lr{\&} \lr{Wagenmakers} (فصل
17) یک \lr{GCM} بیزی و سلسله‌مراتبی می‌سازند.
\end{quote}

\paragraph{خواندن ریاضی}

\begin{itemize}
\tightlist
\item
  \textbf{فاصله} \(\sum_m w_m|x_{im} - x_{em}|\): ویژگی‌هایی را بشمار که
  آیتم جدید \(i\) و مثال ذخیره‌شده‌ی \(e\) در آن‌ها فرق دارند، ولی هر ویژگی
  را با توجه آن یعنی \(w_m\) وزن بده. با 4 ویژگی دودویی و وزن‌های برابر
  0.25، یک ویژگی ناهمخوان فاصله‌ی 0.25 می‌دهد.
\item
  \textbf{شباهت} \(e^{-c \cdot \mathrm{dist}}\): فاصله‌ی 0 شباهت 1 می‌دهد؛
  با \(c = 2\)، فاصله‌ی 0.25 شباهت \(e^{-0.5} = 0.61\) و فاصله‌ی 1 شباهت
  \(e^{-2} = 0.14\) می‌دهد. \(c\) بزرگ‌تر باعث می‌شود شباهت سریع‌تر کم شود،
  پس فقط مثال‌های خیلی نزدیک مهم‌اند.
\item
  \textbf{قاعده‌ی انتخاب:} شباهت‌ها به همه‌ی مثال‌های \lr{A} و همه‌ی مثال‌های
  \lr{B} را جمع بزن و تقسیم کن. در مثال حل‌شده، \lr{T3} شباهت کل 2.79 به
  \lr{A} و 1.09 به \lr{B} دارد، پس \(P(A) = 2.79/3.89 = 0.72\).
\end{itemize}

\subsubsection{\texorpdfstring{3. مدل \lr{drift-diffusion
(DDM)}}{3. مدل drift-diffusion (DDM)}}

\lr{random walk} هفته‌ی 1 را بگیرید و بگذارید اندازه‌ی گام و مدت گام با هم
کوچک شوند. آن‌وقت شواهد \(X(t)\) به‌طور پیوسته تغییر می‌کند:

\[
dX = v\,dt + s\,dW(t), \qquad X(0) = z,
\]

که در آن \(dW\) نویز سفید گاوسی (فرایند \lr{Wiener)} است. وقتی \(X\) به
\(a\) (درست) یا \(0\) (خطا) برسد پاسخ داده می‌شود؛
\(\mathrm{RT} = T_{\text{\lr{decision}}} + T_{er}\).

\begin{longtable}[]{@{}
  >{\raggedright\arraybackslash}p{(\linewidth - 4\tabcolsep) * \real{0.3333}}
  >{\raggedright\arraybackslash}p{(\linewidth - 4\tabcolsep) * \real{0.3333}}
  >{\raggedright\arraybackslash}p{(\linewidth - 4\tabcolsep) * \real{0.3333}}@{}}
\caption{\centering 3. مدل \lr{drift-diffusion (DDM)}}\tabularnewline
\toprule\noalign{}
\begin{minipage}[b]{\linewidth}\raggedright
پارامتر
\end{minipage} & \begin{minipage}[b]{\linewidth}\raggedright
معنا
\end{minipage} & \begin{minipage}[b]{\linewidth}\raggedright
دست‌کاری معمولی که تغییرش می‌دهد
\end{minipage} \\
\midrule\noalign{}
\endfirsthead
\toprule\noalign{}
\begin{minipage}[b]{\linewidth}\raggedright
پارامتر
\end{minipage} & \begin{minipage}[b]{\linewidth}\raggedright
معنا
\end{minipage} & \begin{minipage}[b]{\linewidth}\raggedright
دست‌کاری معمولی که تغییرش می‌دهد
\end{minipage} \\
\midrule\noalign{}
\endhead
\bottomrule\noalign{}
\endlastfoot
\(v\) \lr{(drift rate)} & کیفیت / نرخ انباشت شواهد & دشواری محرک؛ حساسیت
ادراکی؛ توجه \\
\(a\) (فاصله‌ی مرزها) & احتیاط پاسخ (تنظیم سرعت\lr{--}دقت) & دستورالعمل سرعت
در برابر دقت \\
\(z\) (نقطه‌ی شروع) & سوگیری پیشینی به سمت یک پاسخ & نرخ‌های پایه،
پاداش‌ها \\
\(T_{er}\) \lr{(non-decision time}؛ \(t_0\) هم نوشته می‌شود) & زمان
رمزگذاری + حرکتی & نوع پاسخ، تخریب محرک \\
\(s\) (نویز درون \lr{trial)} & ثابت مقیاس & ثابت (طبق قرارداد
\(s = 0.1\) در \lr{EZ} و \(s = 1\) در جاهای دیگر) \\
\end{longtable}

\lr{DDM} کامل تغییرپذیری بین \lr{trial}ها در \lr{drift}، نقطه‌ی شروع و
\lr{non-decision time} را هم دارد. تمرین 5 هفته‌ی 1 نشان داد که
تغییرپذیری \lr{drift} خطاهای کند تولید می‌کند.

\textbf{پیش‌بینی‌های شاخص:} توزیع‌های \lr{RT} چوله به راست؛ \lr{trade-off}
سرعت\lr{--}دقت وقتی \(a\) تغییر می‌کند؛ محرک‌های آسان‌تر که پاسخ‌هایی هم سریع‌تر
\emph{و} هم دقیق‌تر می‌دهند (\(v\) بالاتر)؛ تغییر \(T_{er}\) که کل توزیع
\lr{RT} را بدون تغییر شکلش جابه‌جا می‌کند.

\paragraph{خواندن ریاضی}

\(dX = v\,dt + s\,dW(t)\) این‌طور خوانده می‌شود: «در یک گام زمانی خیلی
کوچک \(dt\)، شواهد به اندازه‌ی یک هل قطعی \(v\,dt\) به‌علاوه‌ی یک ضربه‌ی
تصادفی \(s\,dW\) تغییر می‌کند». \(dW\) نویز گاوسی است که واریانسش متناسب
با \(dt\) است (پس در یک گام به طول \(dt\) ضربه انحراف معیار
\(s\sqrt{dt}\) دارد). جدول پارامترها بخشی است که باید به خاطر سپرد:
\(v\) = کیفیت شواهد، \(a\) = احتیاط، \(z\) = سوگیری، \(T_{er}\) =
\lr{non-decision time.}

\subsubsection{\texorpdfstring{4. \lr{EZ-diffusion}}{4. EZ-diffusion}}

\lr{Wagenmakers}، \lr{van der Maas} و \lr{Grasman} (2007) برآوردهای
بسته‌ی \(v\)، \(a\) و \(T_{er}\) را (با فرض \(z = a/2\)، بدون تغییرپذیری
بین \lr{trial}ها، \(s = 0.1\)) از سه آماره به دست آوردند: نسبت پاسخ‌های
درست \(P_c\)، \textbf{واریانس} \lr{RT}های درست \(V_{RT}\) (بر حسب
\lr{s²)} و \textbf{میانگین} \lr{RT}های درست \(M_{RT}\) (بر حسب \lr{s):}

\[
L = \operatorname{logit}(P_c) = \ln\frac{P_c}{1 - P_c}, \qquad
x = \frac{L\,(L P_c^2 - L P_c + P_c - 0.5)}{V_{RT}},
\]

\[
v = \operatorname{sign}(P_c - 0.5)\; s\; x^{1/4}, \qquad
a = \frac{s^2 L}{v},
\]

\[
y = -\frac{v a}{s^2}, \qquad
\mathrm{MDT} = \frac{a}{2v}\cdot\frac{1 - e^{y}}{1 + e^{y}}, \qquad
T_{er} = M_{RT} - \mathrm{MDT}.
\]

\textbf{حالت‌های مرزی:} فرمول‌ها در \(P_c = 0.5\) (که آن‌وقت \(v = 0\)) و
\(P_c = 1\) (که آن‌وقت \(L = \infty\)) کار نمی‌کنند. برای \(P_c = 1\) یک
راه‌حل استاندارد استفاده از \(P_c = 1 - 1/(2N)\) است، شبیه تصحیح‌های نرخ
حدی هفته‌ی 2.

\textbf{نقاط قوت و محدودیت‌ها.} \lr{EZ} سریع است، به بهینه‌سازی نیاز ندارد
و وقتی فرض‌هایش برقرارند پارامترها را خوب بازیابی می‌کند (بخش \lr{C lab).
RT}های خطا و شکل توزیع \lr{RT} فراتر از میانگین و واریانس را نادیده
می‌گیرد، و نمی‌تواند سوگیری یا تغییرپذیری‌های بین \lr{trial}ها را برآورد
کند.

\paragraph{خواندن ریاضی}

معادله‌های \lr{EZ} ترسناک به نظر می‌رسند ولی یک دستور پخت‌اند؛ لازم نیست
کسی آن‌ها را حفظ کند.

\begin{enumerate}
\def\labelenumi{\arabic{enumi}.}
\tightlist
\item
  \(L = \ln(P_c/(1-P_c))\): \textbf{\lr{logit}} دقت («لگاریتم بخت درست
  بودن»). \(P_c = 0.8\) می‌دهد \(L = \ln 4 = 1.39\).
\item
  \(x\) دقت را با \textbf{واریانس} زمان واکنش پاسخ‌های درست ترکیب می‌کند:
  هرچه زمان‌های واکنش برای یک دقت مشخص پراکنده‌تر باشند، \lr{drift} کمتر
  است.
\item
  بعد \(v\) و \(a\) از \(L\) و \(x\) به دست می‌آیند؛ \(\mathrm{MDT}\)
  \emph{میانگین زمان تصمیم} پیش‌بینی‌شده است.
\item
  \(T_{er} = M_{RT} - \mathrm{MDT}\): هر بخشی از میانگین زمان واکنش که
  زمان تصمیم نیست، \lr{non-decision time} است.
\end{enumerate}

از مثال حل‌شده (بخش 2) برای نشان دادن هر قدم با عدد استفاده کنید؛

\paragraph{برای یک مهندس}

\begin{itemize}
\tightlist
\item
  \textbf{\lr{GCM}} یک طبقه‌بند مبتنی بر شباهت با کرنل نمایی و یک متریک
  \lr{L1} وزن‌دار \lr{(city-block)} است، نزدیک به یک \textbf{طبقه‌بند
  کرنلی / پنجره‌ی \lr{Parzen}}. وزن‌های توجه مثل وزن‌های یادگرفته‌شده‌ی
  ویژگی‌ها عمل می‌کنند؛ \(c\) معکوس پهنای کرنل است.
\item
  \textbf{\lr{DDM}} یک \textbf{فرایند \lr{Wiener} با رانش} بین دو مرز
  جاذب است: مسئله‌ی کلاسیک زمان اولین عبور \lr{(first-passage time).
  SPRT} پیوسته در زمان (کادر هفته‌ی 1) دقیقاً همین فرایند است، با \(v\)
  متناسب با دامنه‌ی سیگنال و \(s\) سطح نویز.
\item
  توزیع‌های زمان واکنش همان \textbf{توزیع‌های زمان اولین عبور} هستند و به
  همان دلیلی چوله به راست‌اند که تأخیرهای آشکارسازی یک آشکارساز ترتیبی.
\item
  \textbf{\lr{EZ-diffusion}} یک تخمین‌گر روش گشتاورها \lr{(method of
  moments)} است: سه گشتاور نمونه (دقت، میانگین و واریانس زمان واکنش‌های
  درست) را با عبارت‌های بسته‌ی مدل برابر می‌گذارد، به جای اینکه درست‌نمایی
  را بیشینه کند.
\end{itemize}

\paragraph{برای مدرس}

\begin{itemize}
\tightlist
\item
  \textbf{بگذارید دانشجوها خودشان شرکت‌کننده باشند.} پیش از \lr{GCM}، 9
  محرک آموزشی ساختار 5-4 را به شکل نقاشی‌های ساده نشان دهید (مثلاً شکل‌هایی
  با 4 ویژگی دودویی)، دسته‌ها را با بازخورد یاد بدهید و بعد \lr{A1} و
  \lr{A2} را آزمون کنید. داده‌ی خود کلاس اغلب الگوی \lr{exemplar} را نشان
  می‌دهد.
\item
  \textbf{\lr{DDM} را با همان تصویر هفته‌ی 1 و یک مسیر هموار و پرنویز
  بکشید.} بعد هر بار یک پارامتر را روی تخته تغییر دهید (مرزهای دورتر،
  \lr{drift} تندتر، شروع دیرتر) و بپرسید دقت و زمان واکنش چه می‌شود.
\item
  \textbf{سؤال‌های رایج.}

  \begin{itemize}
  \tightlist
  \item
    برای دسته‌های بزرگ و پرساختار \lr{prototype} هم خوب کار می‌کند. هر دو
    ایده در مدل‌های امروزی زنده‌اند.
  \item
    \emph{«چرا فقط میانگین زمان واکنش را مقایسه نکنیم؟»} یک گروه کندتر
    ممکن است محتاط‌تر باشد، شواهد بدتری داشته باشد یا از نظر حرکتی کندتر
    باشد؛ فقط \lr{DDM} این‌ها را جدا می‌کند.
  \item
    \emph{«آیا پارامترهای \lr{EZ} و \lr{DDM} کامل یکی‌اند؟»} همان کمیت‌ها
    را برآورد می‌کنند، ولی \lr{EZ} فرض می‌کند سوگیری و تغییرپذیری بین
    \lr{trial}ها وجود ندارد.
  \end{itemize}
\end{itemize}

\begin{center}\rule{0.5\linewidth}{0.5pt}\end{center}

\Needspace{11\baselineskip}

\subsection{مثال حل‌شده}

\subsubsection{\texorpdfstring{بخش 1: یک پیش‌بینی \lr{GCM} با
دست}{بخش 1: یک پیش‌بینی GCM با دست}}

محرک \lr{T3} (1111)، توجه برابر \(w_m = 0.25\)، حساسیت \(c = 2\).

\begin{longtable}[]{@{}llll@{}}
\caption{\centering بخش 1: یک پیش‌بینی \lr{GCM }با دست}\tabularnewline
\toprule\noalign{}
\lr{Exemplar} & ابعاد ناهمخوان & فاصله & شباهت
\(e^{-2\,\mathrm{dist}}\) \\
\midrule\noalign{}
\endfirsthead
\toprule\noalign{}
\lr{Exemplar} & ابعاد ناهمخوان & فاصله & شباهت
\(e^{-2\,\mathrm{dist}}\) \\
\midrule\noalign{}
\endhead
\bottomrule\noalign{}
\endlastfoot
\lr{A1} 1110 & 1 & 0.25 & 0.607 \\
\lr{A2} 1010 & 2 & 0.50 & 0.368 \\
\lr{A3} 1011 & 1 & 0.25 & 0.607 \\
\lr{A4} 1101 & 1 & 0.25 & 0.607 \\
\lr{A5} 0111 & 1 & 0.25 & 0.607 \\
\textbf{جمع \lr{A}} & & & \textbf{2.794} \\
\lr{B1} 1100 & 2 & 0.50 & 0.368 \\
\lr{B2} 0110 & 2 & 0.50 & 0.368 \\
\lr{B3} 0001 & 3 & 0.75 & 0.223 \\
\lr{B4} 0000 & 4 & 1.00 & 0.135 \\
\textbf{جمع \lr{B}} & & & \textbf{1.094} \\
\end{longtable}

\[
P(A\mid \text{\lr{T3}}) = \frac{2.794}{2.794 + 1.094} = 0.719.
\]

\textbf{تفسیر.} \lr{T3} خودِ \lr{prototype} دسته‌ی \lr{A} است، ولی مدل
\lr{exemplar} با این پارامترها فقط حدود 72\% پاسخ \lr{«A»} پیش‌بینی
می‌کند: دو \lr{exemplar} از \lr{B} (یعنی \lr{B1} و \lr{B2)} فقط دو ویژگی
فاصله دارند. یک \(c\) بزرگ‌تر (شباهت تیزتر) این احتمال را بالا می‌برد، چون
نزدیک‌ترین همسایه‌ها (چهار \lr{exemplar} از \lr{A} در فاصله‌ی 0.25) غالب
می‌شوند.

\subsubsection{\texorpdfstring{بخش 2: \lr{EZ-diffusion} با
دست}{بخش 2: EZ-diffusion با دست}}

یک شرکت‌کننده \(P_c = 0.802\)، \(V_{RT} = 0.112\) \lr{s²} و
\(M_{RT} = 0.723\) \lr{s} دارد.

\begin{enumerate}
\def\labelenumi{\arabic{enumi}.}
\tightlist
\item
  \(L = \ln(0.802/0.198) = 1.399\).
\item
  \(x = 1.399 \times (1.399 \times 0.643 - 1.399 \times 0.802 + 0.802 - 0.5)/0.112 = 0.998\).
\item
  \(v = 0.1 \times 0.998^{1/4} = 0.100\).
\item
  \(a = 0.01 \times 1.399/0.100 = 0.140\).
\item
  \(y = -0.100 \times 0.140/0.01 = -1.399\)، پس
  \(\mathrm{MDT} = (0.140/0.200) \times (1 - e^{-1.399})/(1 + e^{-1.399}) = 0.423\)
  \lr{s.}
\item
  \(T_{er} = 0.723 - 0.423 = 0.300\) \lr{s.}
\end{enumerate}

\textbf{تفسیر.} از میانگین \lr{RT} درست 723 میلی‌ثانیه، حدود 423
میلی‌ثانیه صرف انباشت شواهد و 300 میلی‌ثانیه صرف رمزگذاری و پاسخ می‌شود.
\lr{drift} 0.10 و مرز 0.14 (در واحدهای \lr{EZ} با \(s = 0.1\)) مقادیر
معمولی برای یک تکلیف با دشواری متوسط‌اند.

\begin{center}\rule{0.5\linewidth}{0.5pt}\end{center}

\Needspace{11\baselineskip}

\subsection{\texorpdfstring{\lr{lab} (2 ساعت)}{lab (2 ساعت)}}

\subsubsection{اهداف}

\begin{enumerate}
\def\labelenumi{\arabic{enumi}.}
\tightlist
\item
  صورت‌بندی \lr{GCM} و مدل \lr{prototype} برای ساختار 5-4؛ شبیه‌سازی داده
  از \lr{GCM.}
\item
  برازش هر دو مدل با \lr{MLE}، مقایسه‌شان، و بررسی پیش‌بینی کلیدی \lr{A1}
  در برابر \lr{A2.}
\item
  بررسی بازیابی پارامتر برای \lr{GCM.}
\item
  شبیه‌سازی یک فرایند \lr{drift-diffusion}، بازیابی پارامترهایش با
  \lr{EZ-diffusion}، و تفسیر یک آزمایش 2 × 2.
\end{enumerate}

\textbf{زمان‌بندی:} بخش \lr{A} 20 دقیقه · بخش \lr{B} 25 دقیقه · بخش
\lr{C} 15 دقیقه · بخش \lr{D} 40 دقیقه · بخش \lr{E} 20 دقیقه.

\subsubsection{\texorpdfstring{بخش \lr{B:} برازش و مقایسه (25
دقیقه)}{بخش B: برازش و مقایسه (25 دقیقه)}}

\textbf{تفسیر کنید:} 1. (\(\Delta\mathrm{AIC} = 3.3\)، پس وزن‌های
\lr{Akaike} حدود 0.84 در برابر 0.16: ترجیحی متوسط، نه قاطع.) 2. به
\lr{A1} و \lr{A2} نگاه کنید. داده نشان می‌دهد \lr{A2} (0.90) بهتر از
\lr{A1} (0.75) دسته‌بندی شده. \lr{GCM} این ترتیب را پیش‌بینی می‌کند (0.88
در برابر 0.84)؛ مدل \lr{prototype} نمی‌تواند (0.82 در برابر 0.82). چرا
این مقایسه بین دو نظریه تیزتر از آماره‌ی برازش کلی تمایز می‌گذارد؟ 3.
وزن‌های توجه \lr{GCM} را تفسیر کنید.

\subsubsection{\texorpdfstring{بخش \lr{C:} آیا \lr{GCM} در این طرح قابل
بازیابی است؟ (15
دقیقه)}{بخش C: آیا GCM در این طرح قابل بازیابی است؟ (15 دقیقه)}}

\textbf{به چه دقت کنید:} وزن‌های توجه خوب بازیابی می‌شوند \lr{(r} =
0.76\lr{--}0.91)، ولی حساسیت \(c\) در این طرح عملاً \textbf{قابل بازیابی نیست}
\lr{(r} = 0.08). پس باید وزن‌های توجه بخش \lr{B} را تفسیر کنیم ولی مقدار
\(\hat c = 3.82\) را \textbf{نه}. \lr{trial}های بیشتر، دامنه‌ی محرک
گسترده‌تر، یا یک پیشین بیزی روی \(c\) کمک می‌کند. این دقیقاً همان نوع
نتیجه‌ای است که مطالعه‌ی بازیابی برایش ساخته شده (هفته‌ی 7).

\subsubsection{\texorpdfstring{بخش \lr{E:} تفسیر آزمایش 2 × 2 (20
دقیقه)}{بخش E: تفسیر آزمایش 2 × 2 (20 دقیقه)}}

2. توصیف کنید توزیع \lr{RT} از دستورالعمل سرعت به دقت چطور تغییر می‌کند
(نمودار بخش \lr{D).} 3. یک مطالعه‌ی بالینی نشان می‌دهد بیماران از گروه
کنترل کندترند ولی به همان اندازه دقیق.

\subsubsection{تمرین‌ها (به ترتیب سختی)}

\textbf{تمرین 1 \ensuremath{\star}.} \(P(A\mid\text{\lr{A1}})\) و \(P(A\mid\text{\lr{A2}})\) را
برای وزن‌های برابر و \(c = 2\) تحت هر دو مدل با دست حساب کنید. آیا
پیش‌بینی کلیدی را بازتولید می‌کنند؟

\textbf{تمرین 2 \ensuremath{\star}.} \lr{EZ} را روی آماره‌های مثال حل‌شده و روی
\(P_c = 1.0\) با \(N = 100\) \lr{trial} (با تصحیح مرزی) اعمال کنید. بدون
تصحیح چه اتفاقی می‌افتد؟

\textbf{تمرین 3 \ensuremath{\star}\ensuremath{\star} (داده‌ی \lr{prototype).}} داده را از مدل
\textbf{\lr{prototype}} تولید کنید و هر دو مدل را برازش کنید. آیا
مقایسه‌ی مدل، مدل مولد را تشخیص می‌دهد؟ یک بازیابی مدل کوچک اجرا کنید (20
مجموعه داده برای هر مدل مولد) و \lr{confusion matrix} 2 × 2 را بسازید.

\(c \sim \mathrm{Gamma}(2, \text{\lr{scale }} 2)\) ·
\(w \sim \mathrm{Dirichlet}(1,\allowbreak 1,\allowbreak 1,\allowbreak 1)\) پسین وزن‌های توجه را با
\lr{MLE} مقایسه کنید و یک بررسی پیش‌بینی پسین روی تفاوت \lr{A1} و \lr{A2}
اجرا کنید. (برای نسخه‌ی سلسله‌مراتبی، \lr{Lee} \lr{\&} \lr{Wagenmakers} فصل 17
را ببینید.)

\subsubsection{\texorpdfstring{\lr{lab check-off}}{lab check-off}}

نشان دهید: (1) جدول برازش و نمودار پیش‌بینی \lr{GCM} در برابر مدل
\lr{prototype} همراه با جوابتان به سؤال 2، (2) جدول \lr{EZ} برای طرح 2 ×
2، و (3) جوابتان به سؤال 1 بخش \lr{E.}

\begin{center}\rule{0.5\linewidth}{0.5pt}\end{center}

\subsection{\texorpdfstring{ارتباط با هوش مصنوعی و
\lr{LLM}ها}{ارتباط با هوش مصنوعی و LLMها}}

\textbf{مدل‌های \lr{exemplar} در یک «فضای معنایی» یادگرفته‌شده.} \lr{GCM}
به یک فضای روان‌شناختی نیاز دارد تا فاصله‌ها را در آن اندازه بگیرد. به‌طور
سنتی این فضا از رتبه‌بندی‌های شباهت و مقیاس‌بندی چندبعدی \lr{(MDS)} می‌آید.
مدل‌های زبانی یک جایگزین پیشنهاد می‌کنند: هر کلمه یا جمله به یک
\textbf{بردار \lr{embedding}} نگاشت می‌شود، و فاصله‌ی بین بردارها اغلب به
جای شباهت روان‌شناختی استفاده می‌شود. ترکیب \lr{embedding}ها با \lr{GCM}
مدلی می‌دهد که می‌تواند \emph{هر} کلمه‌ی تازه‌ای را بدون رتبه‌بندی شباهت
اضافه دسته‌بندی کند. اینکه فاصله‌های \lr{embedding} با شباهت انسانی جور
درمی‌آیند یا نه یک سؤال تجربی است، و می‌توانید آن را با ابزارهای این درس
بیازمایید.

سه پیش‌بینی چطور عوض می‌شوند، و چرا \lr{«nutcracker»} بیشتر از همه تغییر
می‌کند؟ - \ensuremath{\star}\ensuremath{\star} فرض کنید 20 شرکت‌کننده سه محرک آزمون را دسته‌بندی کرده‌اند، با
19، 1 و 12 پاسخ \lr{«fruit».} \(c\) را با \lr{MLE} برازش کنید (درست‌نمایی
دوجمله‌ای، هفته‌ی 4) و یک آزمون \(G^2\) اجرا کنید. آیا فضای \lr{embedding}
با یک پارامتر برازش خوبی دارد؟ - \ensuremath{\star}\ensuremath{\star} بحث: \lr{embedding}ها از متن یاد
گرفته شده‌اند، نه از ادراک.

\begin{center}\rule{0.5\linewidth}{0.5pt}\end{center}

\subsection{اشتباهات رایج و برداشت‌های نادرست}

\begin{enumerate}
\def\labelenumi{\arabic{enumi}.}
\tightlist
\item
  \textbf{وزن‌های توجه بدون قید.} بدون قید جمع‌برابر-یک، \(c\) و \(w\)
  کاملاً با هم \lr{trade-off} دارند \lr{(non-identifiable).} از
  \lr{reparameterization softmax} استفاده کنید.
\item
  \textbf{قضاوت درباره‌ی مدل‌های \lr{exemplar} و \lr{prototype} فقط با
  برازش کلی.} پیش‌بینی‌هایی که از نظر نظری تشخیصی‌اند \lr{(A1} در برابر
  \lr{A2)} اغلب آگاهی‌بخش‌تر از تفاوت‌های کوچک \lr{AIC} هستند.
\item
  \textbf{میانگین‌گیری داده‌ی دسته‌بندی روی شرکت‌کننده‌هایی} که به ابعاد
  مختلف توجه می‌کنند. این می‌تواند یک فرد \lr{exemplar‌}گونه را
  \lr{prototype‌}گونه نشان دهد (و برعکس)؛ خطاهای میانگین‌گیری هفته‌ی 5 را
  به یاد بیاورید.
\item
  \textbf{استفاده از \lr{EZ} با واحدهای غلط.} \lr{RT}ها باید بر حسب
  \textbf{ثانیه} و واریانس بر حسب \lr{s²} باشد، با \(s = 0.1\). قاطی
  کردن \lr{ms} و \lr{s} نتیجه‌ی بی‌معنی می‌دهد.
\item
  \textbf{استفاده از \lr{EZ} وقتی فرض‌هایش نقض شده‌اند} (سوگیری قوی،
  خطاهای کند یا سریع). \lr{RT}های خطا را بررسی کنید؛ در صورت نیاز از
  \lr{DDM} کامل استفاده کنید.
\item
  \textbf{«گروه کندتر \lr{drift rate} پایین‌تری دارد.»} کندی می‌تواند از
  مرز بالاتر، \lr{drift} پایین‌تر یا \lr{non-decision time} طولانی‌تر
  بیاید.
\item
  \textbf{فراموش کردن حالت‌های مرزی} \(P_c = 1\) یا \(P_c \le 0.5\) در
  \lr{EZ.}
\end{enumerate}

\begin{center}\rule{0.5\linewidth}{0.5pt}\end{center}

\subsection{خودآزمایی}

\textbf{پاسخ} \lr{Exemplar:} یک دسته با همه‌ی نمونه‌های ذخیره‌شده نمایش
داده می‌شود و دسته‌بندی به شباهت جمع‌شده با آن‌ها بستگی دارد.
\lr{Prototype:} یک دسته با یک خلاصه‌ی واحد (گرایش مرکزی) نمایش داده می‌شود
و دسته‌بندی به شباهت با آن خلاصه بستگی دارد.

\textbf{2.} در ساختار 5-4، چرا مقایسه‌ی \lr{A1} و \lr{A2} تشخیصی است؟
\textbf{پاسخ} \lr{A1} به \lr{prototype} دسته‌ی \lr{A} نزدیک‌تر است (3 از 4
ویژگی)، پس مدل \lr{prototype} پیش‌بینی می‌کند بهتر از \lr{A2} دسته‌بندی
شود. ولی \lr{A2} به \lr{exemplar}های دیگر \lr{A} خیلی شبیه است و به هیچ
\lr{exemplar} از \lr{B} شبیه نیست، در حالی که \lr{A1} به دو
\lr{exemplar} از \lr{B} شبیه است، پس مدل \lr{exemplar} می‌تواند \lr{A2} \ensuremath{\geq}
\lr{A1} را پیش‌بینی کند. دو نظریه پیش‌بینی‌های ترتیبی متضاد دارند.

\textbf{3.} پارامترهای \(c\) و \(w_m\) در \lr{GCM} چه چیزی را نشان
می‌دهند؟ \textbf{پاسخ} \(c\) حساسیت است: شباهت با فاصله‌ی روان‌شناختی چقدر
سریع کم می‌شود (\(c\) بزرگ یعنی فقط \lr{exemplar}های خیلی نزدیک مهم‌اند).
\(w_m\) نسبت توجهی است که به بعد \(m\) داده می‌شود و فاصله‌ها را در امتداد
ابعاد مورد توجه کش می‌دهد.

\textbf{پاسخ} (الف) فاصله‌ی مرزها \(a\). (ب) \lr{drift rate} \(v\). (ج)
\lr{non-decision time} \(T_{er}\).

\textbf{پاسخ} نسبت پاسخ‌های درست، واریانس \lr{RT}های درست و میانگین
\lr{RT}های درست. فرض می‌کند نقطه‌ی شروع بی‌سوگیر است (\(z = a/2\))،
پارامترها تغییرپذیری بین \lr{trial}ها ندارند، و مقیاس \(s = 0.1\) است.

\textbf{6.} \lr{DDM} چه ارتباطی با مدل \lr{random walk} هفته‌ی 1 دارد؟
\textbf{پاسخ} حد زمان-پیوسته و شواهد-پیوسته‌ی \lr{random walk} است: وقتی
اندازه‌ی گام و مدت گام با هم به صفر میل کنند، \lr{random walk} به یک
فرایند انتشار \lr{Wiener} با \lr{drift} \(v\) و نویز \(s\) همگرا می‌شود.

\begin{center}\rule{0.5\linewidth}{0.5pt}\end{center}

\subsection{تکلیف}

\begin{itemize}
\tightlist
\item
  \textbf{\lr{A5 (MCMC} از صفر)} در ابتدای درس این هفته تحویل داده
  می‌شود.
\item
  تکلیف جدیدی نداریم: روی \textbf{پروژه‌ی پایانی} کار کنید. به مواد درس
  ارجاع دهید و بگویید چه چیزی را تغییر داده‌اید.
\end{itemize}

\begin{center}\rule{0.5\linewidth}{0.5pt}\end{center}

\subsection{مطالعه‌ی بیشتر}

\begin{itemize}
\tightlist
\item
  \textbf{\lr{Lee} \lr{\&} \lr{Wagenmakers} (2014)، فصل 17.} مدل بافت
  تعمیم‌یافته که به روش بیزی (و سلسله‌مراتبی) روی داده‌ی یادگیری دسته برازش
  می‌شود. \emph{(منبع اصلی برای دسته‌بندی.)}
\item
  \textbf{\lr{Farrell} \lr{\&} \lr{Lewandowsky} (2018)}: فصل مدل‌های زمان پاسخ
  انتخابی {[}شماره‌ی فصل بررسی شود{]} برای مدل \lr{drift-diffusion} و
  خویشاوندانش.
\item
  \textbf{\lr{Verguts} (2022)}: فصل‌های دسته‌بندی و تصمیم‌گیری {[}فصل‌ها
  بررسی شود{]}.
\item
  \emph{اختیاری، خارج از فهرست درس:} \lr{Wagenmakers, E.-J., van der
  Maas, H. L. J.,} \lr{\&} \lr{Grasman, R. P. P. P.} (2007). \lr{An
  EZ-diffusion model for response time and accuracy.}
  \emph{\lr{Psychonomic Bulletin} \lr{\&} \lr{Review,} 14}, 3\lr{--}22. \lr{Medin,
  D. L.,} \lr{\&} \lr{Schaffer, M. M.} (1978). \lr{Context theory of
  classification learning.} \emph{\lr{Psychological Review,} 85},
  207\lr{--}238.
\end{itemize}

\emph{پیش‌نمایش هفته‌ی 15:} یادگیری از پاداش‌ها (یادگیری تقویتی)، ذهن
به‌عنوان یک مشاهده‌گر بیزی، و اینکه مدل‌های شناختی و سیستم‌های هوش مصنوعی
چطور به هم کمک می‌کنند.
